\documentclass[11pt,a4paper]{article}
\usepackage[utf8]{inputenc}
\usepackage[T1]{fontenc}
\usepackage[brazilian]{babel}
\usepackage[margin=2cm]{geometry}
\usepackage{lmodern}
\usepackage{amssymb}
\usepackage{xcolor}
\usepackage{array}
\usepackage{booktabs}
\usepackage{enumitem}
\usepackage{titlesec}
\usepackage{fancyhdr}
\usepackage{parskip}

\definecolor{brandgreen}{HTML}{2E5E3E}
\definecolor{lightbg}{HTML}{F5F2EC}

\titleformat{\section}{\large\bfseries\color{brandgreen}}{}{0pt}{}[{\color{brandgreen}\titlerule}]
\titleformat{\subsection}{\bfseries\color{brandgreen}}{}{0pt}{}

\pagestyle{fancy}
\fancyhf{}
\fancyhead[L]{\small Querência ERP --- Avaliação de usabilidade}
\fancyhead[R]{\small Fluxo 2 --- Diretoria}
\fancyfoot[C]{\small\thepage}
\renewcommand{\headrulewidth}{0.4pt}

\newcommand{\dizer}[1]{%
  \noindent\colorbox{lightbg}{%
    \begin{minipage}{0.95\linewidth}
    \vspace{1mm}
    \textbf{O que dizer ao participante:} ``#1''
    \vspace{1mm}
    \end{minipage}%
  }\vspace{3mm}
}
\newcommand{\criterio}[1]{\noindent\textit{Critério de sucesso objetivo:} #1}
\newcommand{\estado}[1]{\noindent\textit{Estado necessário (já preparado no kit):} #1\vspace{2mm}}

\begin{document}

\begin{center}
  {\LARGE\bfseries\color{brandgreen} Fluxo 2 --- Validação de entrada no evento}\\[1mm]
  {\large RF12, RF20, RF21 --- Papel: Diretoria (painel web)}
\end{center}

\vspace{2mm}
\estado{2 ingressos já pré-emitidos pelo moderador, em nome de associados fictícios: ingresso
\textbf{A} (para ler por câmera) e ingresso \textbf{B} (para digitar o código).}

O QR do ingresso A aparece no celular do moderador (app logado na conta do associado fictício) ou
na impressão do kit; o código do ingresso B está anotado no cartão do kit.

\dizer{Estes dois convidados chegaram ao baile. Confira a entrada de cada um.}

\begin{enumerate}[leftmargin=6mm]
  \item Acessar \textbf{Eventos → [o evento] → Ingressos} (\texttt{/eventos/[id]/ingressos}).
  \item \textbf{Ingresso A:} tocar \textbf{Ler código} (câmera) e apontar para o QR mostrado pelo
    moderador.
  \item \textbf{Ingresso B:} informar o código digitado manualmente no campo de texto, conforme o
    cartão.
  \item Em cada um, na tela de conferência (nome + código do ingresso, evento, categoria, valor,
    status), revisar os dados e tocar \textbf{Confirmar entrada} --- o check-in não é automático
    ao ler o QR.
  \item Confirmar que a listagem/contadores do evento (emitidos / validados / aguardando)
    atualizam.
\end{enumerate}

\textbf{Verificação objetiva:}
\begin{small}
\begin{verbatim}
node scripts/verificar-tarefas.js 2 <CPF-DO-ASSOCIADO-FICTICIO-A> \
  "<trecho do nome do evento>"
node scripts/verificar-tarefas.js 2 <CPF-DO-ASSOCIADO-FICTICIO-B> \
  "<trecho do nome do evento>"
\end{verbatim}
\end{small}
\criterio{ingresso emitido \textbf{e} entrada validada, com horário preenchido, nos dois.}

\vspace{2mm}
\noindent\textit{Se a câmera falhar (permissão, luz, rede), registrar o motivo e pedir o ingresso
A pelo caminho manual, contando como intervenção --- sem usar o ingresso B (que continua reservado
para o caminho manual).}

\vfill
\noindent\footnotesize Se a tarefa travar por mais de \textasciitilde2 min, aplicar 1 intervenção
pontual e registrar; se mesmo assim não avançar, parar e marcar como \textit{não concluída} ---
protocolo completo no \texttt{roteiro-moderador.pdf}.

\end{document}
